\documentclass[5pt, a4paper]{article}

\usepackage[utf8]{inputenc}
\usepackage[T1]{fontenc}
\usepackage{textcomp}
\usepackage{amsmath, amssymb}
\usepackage{multicol}
\usepackage{proof}
\usepackage{enumitem}
\usepackage{fancyhdr}
\usepackage[margin=1in]{geometry}
% figure support
\usepackage{import}
\usepackage{xifthen}
\pdfminorversion=7

\usepackage[most]{tcolorbox}
\usepackage{pdfpages}
\usepackage{transparent}
\newcommand{\incfig}[1]{%
    \def\svgwidth{\columnwidth}
    \import{./figures/}{#1.pdf_tex}
}

\tcbset{
colback=gray!5,
colframe=blue!50!white,
coltitle=black,
arc=0pt,
boxrule=-0.1mm,
fonttitle=\bfseries,
coltitle=black,
left=2.5mm,
leftrule=1mm,
right=3.5mm,
colbacktitle=black!10!white,
toptitle=0.75mm,
bottomtitle=0.5mm,
}
\begin{document}
\begin{tcolorbox}[title=Question 1]
    Use the Intermediate Value theorem to find an interval of length one that
    contains a root of the equation, 
    \begin{enumerate}
        \item $x^3 = 9$
        \item $\ln x + x^2 = 3$
    \end{enumerate}
\end{tcolorbox} 
\begin{tcolorbox}[title=Question 2]
\end{tcolorbox} 
\begin{tcolorbox}[title=Question 3]
\end{tcolorbox} 
\begin{tcolorbox}[title=Question 4]
\end{tcolorbox} 
\begin{tcolorbox}[title=Question 5]
\end{tcolorbox} 
\begin{tcolorbox}[title=Question 6]
\end{tcolorbox} 
\begin{tcolorbox}[title=Question 7]
\end{tcolorbox} 
\end{document}
